% ===== PRÉAMBULE CANONIQUE — à coller VERBATIM en tête de CHAQUE .tex =====
\documentclass[11pt,a4paper]{article}
\usepackage[utf8]{inputenc}\usepackage[T1]{fontenc}\usepackage{lmodern}
\usepackage[french]{babel}
\usepackage{amsmath,amssymb,mathtools}\usepackage{icomma}
\usepackage[version=4]{mhchem}          % \ce{...} : équations & formules
\usepackage{siunitx}\sisetup{locale=FR} % \qty{}{}, \si{}, \ang{}
\usepackage{chemfig}                    % \chemfig{...} : structures développées
\usepackage{xcolor}\usepackage{colortbl}
\usepackage{tikz}\usepackage{pgfplots}\pgfplotsset{compat=1.18}
\usepackage[most]{tcolorbox}\usepackage{enumitem}\usepackage{array,booktabs}
\usepackage[a4paper,margin=2.2cm]{geometry}
\usetikzlibrary{arrows.meta,calc,angles,quotes,positioning,patterns,backgrounds,shapes.geometric}
\definecolor{primary}{RGB}{222,49,99}\definecolor{primarydark}{RGB}{150,22,55}
\definecolor{defcol}{RGB}{0,114,178}\definecolor{methcol}{RGB}{52,92,148}
\definecolor{excol}{RGB}{0,135,90}\definecolor{attncol}{RGB}{200,40,55}
\tcbset{boxbase/.style={enhanced,breakable,boxrule=0.9pt,arc=2.5pt,
  left=3mm,right=3mm,top=2.3mm,bottom=2.3mm,fonttitle=\bfseries,coltitle=white}}
% --- boîtes TUTORIEL (noms EXACTS, ne pas renommer) ---
\newtcolorbox{cle}[1][]{boxbase,colback=defcol!6,colframe=defcol,
  title={\textsc{À retenir}\ifx\relax#1\relax\relax\else\textmd{ : \itshape #1}\fi}}
\newtcolorbox{methode}[1][]{boxbase,colback=methcol!7,colframe=methcol,sharp corners,
  title={\textsc{Méthode}\ifx\relax#1\relax\relax\else\textmd{ --- \itshape #1}\fi}}
\newtcolorbox{exemplebox}[1][]{boxbase,colback=excol!5,colframe=excol,
  title={\textsc{Exemple}\ifx\relax#1\relax\relax\else\textmd{ : \itshape #1}\fi}}
\newtcolorbox{piege}[1][]{boxbase,colback=attncol!6,colframe=attncol,
  title={\textsc{Piège}\ifx\relax#1\relax\relax\else\textmd{ : \itshape #1}\fi}}
\newtcolorbox{remarque}[1][]{boxbase,colback=black!4,colframe=black!45,
  title={\textsc{Remarque}\ifx\relax#1\relax\relax\else\textmd{ : \itshape #1}\fi}}
% --- boîtes EXERCICE / CORRIGÉ (noms EXACTS, auto-comptées) ---
\newtcolorbox[auto counter]{exo}[1][]{boxbase,colback=primary!4,colframe=primary,
  title={\textsc{Exercice}~\thetcbcounter\ifx\relax#1\relax\relax\else\textmd{ --- \itshape #1}\fi}}
\newtcolorbox[auto counter]{corrige}[1][]{boxbase,colback=excol!4,colframe=excol,
  title={\textsc{Corrigé}~\thetcbcounter\ifx\relax#1\relax\relax\else\textmd{ --- \itshape #1}\fi}}
\newcommand{\R}{\mathbb{R}}\newcommand{\N}{\mathbb{N}}\newcommand{\Z}{\mathbb{Z}}\newcommand{\ds}{\displaystyle}
% (un \usepackage{fancyhdr}+en-tête est possible ; il est ignoré côté web)
% =========================================================================
\begin{document}
\begin{corrige}[l'ion formé et son gaz noble]
Métaux : perdent leurs électrons de valence ; non-métaux : en gagnent.
\begin{enumerate}[label=\alph*)]
  \item $\ce{O}$ ($6$ val.) gagne $2$ e$^-$ $\to \ce{O^{2-}}$, $10$ électrons : configuration de $\ce{Ne}$.
  \item $\ce{Na}$ ($1$ val.) perd $1$ e$^-$ $\to \ce{Na^{+}}$, $10$ électrons : configuration de $\ce{Ne}$.
  \item $\ce{Mg}$ ($2$ val.) perd $2$ e$^-$ $\to \ce{Mg^{2+}}$, $10$ électrons : configuration de $\ce{Ne}$.
  \item $\ce{K}$ ($1$ val.) perd $1$ e$^-$ $\to \ce{K^{+}}$, $18$ électrons : configuration de $\ce{Ar}$.
\end{enumerate}
\end{corrige}

\begin{corrige}[quels ions ont la configuration d'un gaz noble ?]
$n(e^-) = Z - q$ ; on compare aux nombres $2, 10, 18\dots$
\begin{enumerate}[label=\alph*)]
  \item $\ce{F^{-}}$ : $9+1 = 10$ électrons \textbf{\textcolor{excol}{$\to$ oui}}, configuration de $\ce{Ne}$.
  \item $\ce{Na^{+}}$ : $11-1 = 10$ \textbf{\textcolor{excol}{$\to$ oui}}, configuration de $\ce{Ne}$.
  \item $\ce{Ca^{2+}}$ : $20-2 = 18$ \textbf{\textcolor{excol}{$\to$ oui}}, configuration de $\ce{Ar}$.
  \item $\ce{H^{+}}$ : $1-1 = 0$ électron \textbf{\textcolor{attncol}{$\to$ non}} : ce n'est la
        configuration d'aucun gaz noble.
\end{enumerate}
\end{corrige}

\begin{corrige}[famille à partir du numéro atomique]
\begin{enumerate}[label=\alph*)]
  \item $Z=11$ : $(\mathrm{K})^2(\mathrm{L})^8(\mathrm{M})^1$, $1$ e$^-$ de valence : \textbf{alcalin} ($\ce{Na}$).
  \item $Z=17$ : $(\mathrm{K})^2(\mathrm{L})^8(\mathrm{M})^7$, $7$ e$^-$ de valence : \textbf{halogène} ($\ce{Cl}$).
  \item $Z=18$ : $(\mathrm{K})^2(\mathrm{L})^8(\mathrm{M})^8$, $8$ e$^-$ de valence : \textbf{gaz noble} ($\ce{Ar}$).
  \item $Z=20$ : $(\mathrm{K})^2(\mathrm{L})^8(\mathrm{M})^8(\mathrm{N})^2$, $2$ e$^-$ de valence :
        \textbf{alcalino-terreux} ($\ce{Ca}$).
\end{enumerate}
\end{corrige}

\begin{corrige}[classer par électronégativité]
L'électronégativité croît vers la droite d'une période.
\begin{enumerate}[label=\alph*)]
  \item $\ce{Na} < \ce{Al} < \ce{P} < \ce{Cl}$.
  \item $\ce{Li} < \ce{N} < \ce{F}$.
\end{enumerate}
\end{corrige}

\begin{corrige}[classer par rayon atomique]
Le rayon croît vers le bas d'une colonne et vers la gauche d'une période.
\begin{enumerate}[label=\alph*)]
  \item $\ce{Li} < \ce{Na} < \ce{K}$ (il augmente en descendant la colonne).
  \item $\ce{Cl} < \ce{Mg} < \ce{Na}$ (il augmente en allant vers la gauche).
\end{enumerate}
\end{corrige}

\end{document}
